\begin{abstract}
Small treewidth alone does not make nonconvex mixed-integer optimization
tractable: box-constrained nonconvex quadratic programming is strongly
NP-hard at treewidth two. We show that two further quantities make tree
decompositions useful for certified global optimization: upper bounds $L_i$
on the curvature of the objective $F$ along each coordinate, and quadratic
growth. Subtracting $L_iw^2/8$ at each grid node whose longest adjacent grid
interval has length $w$ turns exact dynamic programming over a product grid
into a valid lower bound. Min-marginals then filter the grids, and the stage
grids form a certificate, checked by recomputation, that needs no growth
assumption. Under growth in the $L_i$-weighted norm, graded grids keep
$O(\sqrt{\bar\kappa}\log(n+2))$ nodes per coordinate at every accuracy,
where $\bar\kappa$ is a scale-invariant ratio of curvature to growth. For
mixed-integer box quadratic programs with bag size $p$ and input length $I$,
this certifies accuracy $2^{-q}$ in $f(p,\bar\kappa)(I+q+1)^5$ bit
operations; under growth at a unique minimizer $x^*$, with the Euclidean
ratio $\kappa\ge\bar\kappa$, computing an exact minimizer takes
$f_1(p,\kappa)I^{O(1)}$ bit operations. The growth constant is never
needed. Growth with ratio $\kappa$ forces the Hessian block of the interior continuous
coordinates to have smallest eigenvalue at least $2\max_iL_i/\kappa$, and
every other local minimizer $x'$ to satisfy
$F(x')-\min F\ge\max_iL_i\norm{x'-x^*}^2/\kappa$. Unless P${}={}$NP,
neither parameter can be dropped and the dependence on $\kappa$ cannot be polylogarithmic;
under the randomized exponential-time hypothesis, the exponent of $\kappa$
in a bound $f(p)\kappa^{a(p)}I^{O(1)}$ cannot be $o(p)$. We extend the
method to exact partial minimization, totally unimodular coupling
constraints and several minimizers, and test an exact-arithmetic
implementation.
\par\medskip\noindent\textbf{Keywords:} global optimization; mixed-integer quadratic
programming; tree decomposition; dynamic programming; quadratic growth; certificates;
parameterized complexity
\par\noindent\textbf{Mathematics Subject Classification (2020):} 90C26, 90C11, 90C20, 90C39,
90C60, 68Q27
\end{abstract}
